\section{Joint-space control}
\subsection{Joint-space velocity control}
\subsubsection{Point to point control}
For point to point motion,
there is P-, PD- and PI-controller with feedback on positions.
For example,

\begin{algorithm}[H]
\caption{moveJP}
\begin{algorithmic}[1] % The [1] adds line numbers
\Procedure{moveJP}{cfg, robot\_manager, q\_desired}
\While{not breakFlag}
\State $ e \gets q_{desired} - q$ 
    \State $v_{cmd} = K_p \cdot e$
    \If {$e < $ goal\_error}
    \State breakFlag $\gets$ True
    \EndIf
\EndWhile
\EndProcedure
\end{algorithmic}
\label{cart-reference-algorithm}
\end{algorithm}

As one might expect, the P-controller struggles to converge,
while the PI-controller can overshoot.
PD control seems to do best. None of them are \textit{great}.
The robot moves more smoothly and happily if provided
with a nice trajectory. This is the case
because then the errors are smaller. Thankfully,
generating a nice trajectory in joint space is very easy ---
simply follow the straight line from $q$ to $q_\text{desired}$
in joint space, and put a polynomial time law on it.
But first, we introduce trajectory following control.



\subsubsection{Trajectory tracking control}
The control of choise is a P-controller with feedback on positions 
and feedforward velocity term. This effectively forms PD-control.

\begin{algorithm}[H]
\caption{followParametrizedJointTrajectory}
\begin{algorithmic}[1] % The [1] adds line numbers
	\Procedure{followParametrizedJointTrajectory}{cfg, robot\_manager, $q_{ref}$}
\While{not breakFlag}
\State $v_{cmd} \gets K_p (q_{ref}(t) - q) + \dot{q}_{ref}(t)$
\If {$t > t_{final}$ and $ \left\lVert v_{ cmd } \right \rVert < \epsilon   $}
\State breakFlag $\gets$ True
\EndIf
\EndWhile
\EndProcedure
\end{algorithmic}
\label{algo-joint-traj-tracking}
\end{algorithm}

\subsubsection{Most common joint space controller}
The most common joint space task is a point to point motion
to the robot's home position, or otherwise get the robot
close to the relevant part of the environment. 
By starting from a known home position
one can ensure good performance of subsequent controllers.
This is also good practise to ensure repeatability.
However, as mentioned, the point-to-point controllers are not great.
However$^2$, as advertised, it is very easy to find a nice smooth
trajectory for joint-space motions,
and then apply the trajectory tracking controller.

Namely, the smooth path is a straight line in joint space,
$q_{ref}(s) = q_0 + s(q_{desired} - q_0), s \in [0,1]$.
To ensure smooth robot motion, we construct an appropriate 
trajectory. Polynomials are a great and easy solution,
as they are fully determined by their start and end points.
The higher the degree of polynomial, the more descriptive
the conditions can be (setting not only velocity, but also
acceleration, jerk,...). Furthermore, the higher the polynomial,
the more trajectory derivatives are continuous.
Usually, we care about accelerations and use a quintic polynomials
(as then the acceleration reference is a nice cubic instead of
being continuous or linear).
But, for simplicity, here we use a cubic polynomial,
which then has a quadratic velocity profile and a linear acceleration
profile.


\begin{algorithm}[H]
\caption{moveJWTraj}
\begin{algorithmic}[1] % The [1] adds line numbers
\Procedure{moveJWTraj}{cfg, robot\_manager, $q_{desired}$}
\State $ v_{\max} \gets v[argmax(|q_{desired} - q |)] $
\State $T_{final} = 0.8\frac{max(|q_{desired} - q |)}{v_{\max}} $
\State $q_{ref}(t) \gets q_0 + \left( \frac{3t^2}{T_{final}^2} - \frac{2t^3}{T_{final}^3} \right) (q_{desired} - q_0)$
\State $\dot{q}_{ref}(t) \gets \left( \frac{6t}{T_{final}^2} - \frac{6t^2}{T_{final}^3} \right) (q_{desired} - q_0)$
\State followParametrizedJointTrajectory(cfg, robot\_manager, $q_{ref}$)
\EndProcedure
\end{algorithmic}
\label{algo-movejwtraj}
\end{algorithm}

\subsection{Joint-space torque control}
TODO
